% Wahrscheinlichkeit – bank/wahrscheinlichkeit/ (zusammenbau v0.4, Heft)
\einheitenkopf[t8]{Wahrscheinlichkeit}
\setcounter{aufgabe}{4}

\begin{aufgabe}{Zählen – Prüfungsaufgaben}
\begin{teile}
\teil Zwei gleiche Spielsteine mit einer roten und einer blauen Seite werden gleichzeitig geworfen. Wie viele verschiedene Ergebnisse sind möglich? \hfill (P10 2017 OS)
\end{teile}
\end{aufgabe}

\begin{aufgabe}{Einstufig – Prüfungsaufgaben}
\begin{teile}
\teil Bei einer Tombola liegen 45 Nieten und 15 Gewinnlose in der Trommel. Wie groß ist die Wahrscheinlichkeit für einen Gewinn? \hfill (P10 2014 OS)
\teil Eine Kiste enthält 60 Eier, 4 davon sind zerbrochen. Ein Ei wird zufällig genommen. Wie groß ist die Wahrscheinlichkeit, dass es zerbrochen ist? \hfill (P10 2016 OS)
\teil Ein Glücksrad hat acht gleich große Felder mit den Zahlen 1 bis 8. Wie groß ist die Wahrscheinlichkeit für weder 1 noch 8? \hfill (P10 2014 OS)
\teil Ein Glücksrad hat 4 gleich große Felder in Rot und Gelb. Die Wahrscheinlichkeit für Gelb ist 75\,\%. Wie viele Felder sind gelb? \hfill (P10 2018 OS)
\teil Ein Glücksrad hat gleich große Felder; 6 davon sind rot, die übrigen weiß. Die Wahrscheinlichkeit für Rot soll $\frac{3}{4}$ sein. Zeichne das ganze Glücksrad mit den fehlenden weißen Feldern. \hfill (P10 2019 OS)

\kreisleer[2]
\end{teile}
\end{aufgabe}

\begin{aufgabe}{Topf mit P = 50 \% auswählen – Prüfungsaufgaben}
\begin{teile}
\teil Drei Dosen enthalten je 8 Bonbons: links 5 Kirsch und 3 Zitrone, in der Mitte 4 Kirsch und 4 Zitrone, rechts 2 Kirsch und 6 Zitrone. Aus welcher Dose zieht man Kirsch mit der Wahrscheinlichkeit 50\,\%? Kreuze an. \hfill (P10 2015 OS) \\ \kreuz{linke Dose} \\ \kreuz{mittlere Dose} \\ \kreuz{rechte Dose}
\end{teile}
\end{aufgabe}
